\documentclass[12pt]{article}
\usepackage[utf8]{inputenc}
\usepackage[T1]{fontenc}
\usepackage{hyperref}
\title{Sheni Language Resources for K-12 Curriculum Development}
\author{Compiled by Claude Sonnet 5.5 (claude-sonnet-5-5) for Andr\'as Kornai}
\date{}
\begin{document}
\maketitle

\section{Introduction}
Sheni (ISO 639-3 scv; also Shani, Seni, Shenanci) is an East Kainji language of Kaduna State, Nigeria. It is closely tied to Ziriya, which is listed for Bauchi State, and the speakers regard the two as the same language. A 2004 survey found six fluent speakers and perhaps 10--15 semi-fluent speakers. The community has shifted to Hausa, calls itself Shenawa and calls the language Shenanci. According to the sources found, revival is not a priority for the community and the loss of the language is seen as a fait accompli. Do not confuse it with the Chadic ``Shani'' people of Borno and Adamawa, which is a different group. The searches found no school materials, only academic survey records.

\section{Primary Educational Resources}
None found.

\section{Linguistic and Reference Materials}
\begin{itemize}
\item Wikipedia, Sheni-Ziriya language. \url{https://en.wikipedia.org/wiki/Sheni-Ziriya_language}
\item Endangered Languages Project entry for Sheni. \url{https://endangeredlanguages.com/lang/490}
\item Bibliography for that entry. \url{https://endangeredlanguages.com/lang/490/bibliography?hl=de}
\item Wikipedia, East Kainji languages. \url{https://en.wikipedia.org/wiki/East_Kainji_languages}
\item ELDP archive, Documentation of Kainji languages (a 2011 McGill and Blench survey). The search summary mentioned wordlists for a language listed as ``Shen'', but I did not confirm that this is Sheni. \url{https://eldp.access.preservica.com/dk0145/}
\end{itemize}

\section{Religious and Community Resources}
None found.

\section{Digital and Audio Resources}
None confirmed. The archive above may hold wordlists and recordings, but access conditions were not checked.

\section{Curriculum Development Context}
With so few speakers, and community reports of no interest in revival, a school curriculum is not realistic now. If the community's view changes, the first step would be archival work on the wordlists and a consultation with the Shenawa, probably as a heritage-culture module taught in Hausa or English. Do not present any content as language instruction before it has been validated with the remaining speakers.
\end{document}